\documentclass[a4paper, 12pt]{article}
\usepackage[utf8]{inputenc}
\usepackage[T1]{fontenc}
\usepackage[slovak]{babel}
\usepackage{geometry}
\usepackage{graphicx}
\usepackage{url}
\usepackage{listings}
\usepackage{xcolor}
\usepackage{amsmath}
\usepackage{beramono}
\usepackage{float}

\definecolor{vscode_bg}{RGB}{30,30,30}
\definecolor{light_bg}{RGB}{245,245,247}
\definecolor{codeblue}{RGB}{0,110,180}
\definecolor{codegreen}{RGB}{34,139,34}
\definecolor{codepurple}{RGB}{128,0,128}

\lstset{
    language=Python,
    backgroundcolor=\color{light_bg},
    commentstyle=\color{codegreen},
    keywordstyle=\color{codeblue}\bfseries,
    numberstyle=\tiny\color{gray},
    stringstyle=\color{orange},
    basicstyle=\ttfamily\footnotesize,
    columns=flexible,
    fontadjust=true,
    breaklines=true,
    captionpos=b,
    keepspaces=true,
    numbers=left,
    numbersep=10pt,
    showstringspaces=false,
    frame=lines,
    rulecolor=\color{gray!10},
    xleftmargin=15pt
}

\geometry{a4paper, left=2.5cm, right=2.5cm, top=2.5cm, bottom=2.5cm}

\title{Zadanie 2: Kalibrácia kamery a spracovanie obrazu}
\author{Michal Balogh \\ Jana Kampová}
\date{Predmet: PVSO \\ 2025/2026}

\begin{document}

\maketitle
\clearpage
\tableofcontents
\clearpage
\section{Kalibrácia kamery (Task 1)}

Proces kalibrácie slúži na odstránenie optických chýb šošovky. Celý postup prebieha v nasledujúcich krokoch:

\subsection{Príprava súradníc}
Pri kalibrácii pracujeme s dvoma sadami bodov, ktoré sa algoritmus snaží spárovať:

\begin{itemize}
    \item \textbf{Ideálne súradnice (\textbf{objpoints})}: Predstavujú „3D body v reálnom svete“ (ideálnu mriežku šachovnice).
    \begin{itemize}
        \item \textbf{\textbf{objp}}: Pole nul pre $x, y, z$ súradnice, kde $z$ ostáva 0 a $x, y$ sa na seba párujú podľa počtu rohov šachovnice.
    \end{itemize}
    \item \textbf{Obrazové súradnice (\textbf{imgpoints})}: „2D body v rovine obrazu“, teda súradnice z reálneho obrázka šachovnice.
\end{itemize}

\subsection{Detekcia}

Pre každú  snímku hľadáme vnútorné rohy pomocou \textbf{cv2.findChessboardCorners}.
Ak sa nájdu všetky rohy (\textbf{ret == True}), nasleduje :

\begin{itemize}
    \item \textbf{\textbf{cv2.cornerSubPix}}: Táto funkcia hľadá najpresnejšiu súradnicu hrany.
    \item \textbf{Princíp symetrie}: Algoritmus analyzuje okno $11 \times 11$ pixelov a zisťuje, kde je matematický gradient najviac symetrický podľa jasu. Tento bod sa považuje za presný matematický stred rohu (+).
    \item \textbf{Vizuálna kontrola}: Pomocou \textbf{cv2.drawChessboardCorners} overujeme, či program našiel rohy v správnom poradí (od červenej po modrú farbu).
\end{itemize}
\begin{figure}[h]
    \centering
    \includegraphics[width=0.5\textwidth]{chessboard.png}
    \caption{Ukážka detekcie vnútorných rohov na šachovnici v našom programe pomocou \textbf{cv2.drawChessboardCorners}.}
\end{figure}
\subsection{Výpočet modelu a geometria kamery}
Samotný výpočet \textbf{cv2.calibrateCamera} porovnáva ideálnu mriežku s pokrivenými bodmi z fotiek.

\subsubsection*{Vnútorné parametre (Camera Matrix $K$)}
Matica $3 \times 3$ matematicky prepojuje 3D priestor s 2D fotkou.
\[
K = \begin{bmatrix}
f_x & 0 & c_x \\
0 & f_y & c_y \\
0 & 0 & 1
\end{bmatrix}
\]
V matici sú kľúčové parametre $f_x, f_y$ (ohnisková vzdialenosť) a $c_x, c_y$ (optický stred).

\subsubsection*{Model skreslenia (\textbf{dist})}
Reálne šošovky mávajú geometrické deformácie, ktoré v kóde popisujeme koeficientami $(k_1, k_2, p_1, p_2, k_3)$. Rozlišujeme nasledujúce typy:


\begin{itemize}
    \item \textbf{Radiálne skreslenie ($k_n$)}: Spôsobené tvarom šošovky. Spôsobuje ohýbanie priamok smerom od stredu alebo k stredu obrazu.
    \item \textbf{Tangenciálne skreslenie ($p_n$)}: Vzniká pri výrobnej nepresnosti, kedy šošovka nie je dokonale rovnobežná so snímačom (Decentering distortion).
\end{itemize}

\begin{figure}[H]
    \centering
    \includegraphics[width=0.3\textwidth]{image_1b9dfe.png} \hfill
    \includegraphics[width=0.3\textwidth]{image_1b9e1b.png}
    \hfill
    \includegraphics[width=0.3\textwidth]{image_1b9b52.png}
    \hfill
    \includegraphics[width=0.3\textwidth]{image_1b9e3e.png}

    \caption{Zľava: radiálne skreslenie typu Barrel a Pincushion, tangenciálne skreslenie, ideálna mriežka bez skreslenia}
\end{figure}

\subsection{Odstránenie skreslenia (Undistortion)}
Pre narovnanie obrazu používame dve metódy definované v \textbf{UndistortMethod}:

\begin{enumerate}
    \item \textbf{Metóda CROP}: \newline Počíta sa pre každý obrazok. Algoritmus vezme súradnicu rovného pixelu, vypočíta, kde by bol v pokrivenom obraze, a pomocou \textbf{interpolácie} (priemer zo 4 susedných pixelov) mu priradi farbu.
    \item \textbf{Metóda REMAPP}: \newline Najskôr vytvorí mapu rozdielov pre každý jeden pixel. Táto mapa sa vytvorí len raz na začiatku a potom sa už len aplikuje na každý nasnímaný obrázok, čo je dobre využiteľné práve pre spracovanie videa v reálnom čase.
\end{enumerate}

\subsubsection*{Optimalizácia a ROI}
Pomocou \textbf{cv2.getOptimalNewCameraMatrix} a parametra \textbf{alpha=1} (ponechanie všetkých pixelov) získavame \textbf{ROI} (Region of Interest). \newline ROI sú súradnice čistého obdĺžnika bez čiernych rohov $[x, y, rw, rh]$.
Vďaka tomu môžeme urobiť finálny orez (\textbf{dst[y:y+rh, x:x+rw]}), čím zahodíme prázdne čierne rohy, ktoré vznikli narovnaním.

\subsection{Analýza presnosti (Mean Reprojection Error)}
Spätná kontrola prebieha tak, že vezmeme 3D body ideálnej šachovnice a cez našu novú maticu ich premietneme naspäť na fotku.
\begin{itemize}
    \item Používame normu \textbf{NORM\_L2} (Pytagorova veta), ktorá vypočíta priamu vzdialenosť medzi vypočítaným a skutočným bodom.
    \item Dosiahnutý výsledok v pixeloch určuje kvalitu: $0.1 - 0.5$ px je veľmi dobrý výsledok pre bežnú kameru.
\end{itemize}
\subsection{Získané parametre kamery}
Na základe spracovania sady kalibračných obrázkov sme získali nasledujúce parametre :

\begin{itemize}
    \item \textbf{Mean Reprojection Error} = 0.245 px
    \item \textbf{Matica kamery $K$:}
    \[
    K = \begin{bmatrix}
    645.12 & 0 & 320.55 \\
    0 & 644.89 & 240.12 \\
    0 & 0 & 1
    \end{bmatrix}
    \]
    \item \textbf{Koeficienty skreslenia:}
    \textbf{dist} = $[-0.28, 0.09, 0.001, 0.002, -0.01]$
\end{itemize}
Tieto hodnoty potvrdzujú mierne súdkovité skreslenie, ktoré je typické pre širokouhlé webkamery.

\newpage
\newpage
\section{Farebný filter (Task 3)}

\subsection{Farebný priestor HSV}
Pre detekciu farieb sme použili model \textbf{HSV}. Hlavným dôvodom je, že na rozdiel od BGR, HSV oddeľuje jasovú zložku (Value) od farebnej (Hue). To nám umožňuje nastaviť filter, ktorý je odolný voči tieňom a zmenám osvetlenia v miestnosti.


\begin{figure}[h]
    \centering
    \includegraphics[width=0.6\textwidth]{colors.png}
    \caption{HSV farebný model používaný aj v knižnici OpenCV.}
    \label{fig:hsv_wheel}
\end{figure}
\subsubsection*{Škálovanie hodnôt v OpenCV}
Na rozdiel od bežného matematického modelu, OpenCV upravuje rozsahy hodnôt pre potreby spracovania obrazu:
\begin{itemize}
    \item \textbf{Odtieň (H):} Štandardný rozsah $0^{\circ}$--$360^{\circ}$ je vydelený dvoma, čím vzniká rozsah \textbf{0--180}.
    \item \textbf{Sýtosť (S) a Jas (V):} Pôvodný rozsah $0$--$100\%$ je prepočítaný na hodnoty \textbf{0--255}.
\end{itemize}
 V kóde sme pre jednotlivé farby určili nasledujúce hodnoty. \newline Červená farba vyžaduje spojenie dvoch masiek lebo v HSV kruhu sa nachádza na oboch koncoch intervalu (pri 0° aj 180°).

\begin{table}[h]
\centering
\begin{tabular}{|l|c|c|c|}
\hline
\textbf{Farba} & \textbf{Hue (Odtieň)} & \textbf{Saturation (Sýtosť)} & \textbf{Value (Jas)} \\ \hline
Červená (1)    & 0 -- 10               & 100 -- 255                  & 100 -- 255          \\ \hline
Červená (2)    & 170 -- 180            & 100 -- 255                  & 100 -- 255          \\ \hline
Zelená         & 50 -- 70              & 50 -- 255                   & 50 -- 255           \\ \hline
Modrá          & 110 -- 130            & 50 -- 255                   & 50 -- 255           \\ \hline
\end{tabular}

\caption{HSV rozsahy}
\end{table}
\newpage
\subsection{Implementácia}

\begin{lstlisting}[caption={Implementácia farebného filtra pre 3 a 4-kanálové kamery}]
# Nastavenie cielovej farby pre modry filter

if state == SwapState.BLUE:

    # Hladame modru (100-140), ignorujeme tiene (V>50) a bledo-sive odtiene (S>50)

    maska = cv2.inRange(hsv, (100, 50, 50), (140, 255, 255))
    nova_farba_bgr = (0, 255, 255) # BGR kod pre zltu

# Aplikacia podla typu pripojeneho zariadenia
if num_channels == 3:  # Notebook (BGR)
    camera = ()
    image[maska > 0] = nova_farba_bgr + camera

elif num_channels == 4: # Ximea (BGRA)
    camera = (255,)
    image[maska > 0] = nova_farba_bgr + camera


\end{lstlisting}
\subsection{Výsledky filtrácie}
Na nasledujúcom obrázku je zobrazená funkčnosť filtra pre jednotlivé farebné kanály. Pôvodný obraz je porovnaný s výsledkom po aplikácii HSV masky, kde detekovaná červená farba je nahradená ružovou, detekovaná modrá je nahradená žltou a detekovaná zelená je nahradená oranžovou.

\begin{figure}[h]
    \centering
    \includegraphics[width=0.3\textwidth]{pastelk.png}
    \caption{Ukážka ako funguje naša implementácia}
    \label{fig:pencils}
\end{figure}
\newpage

\begin{thebibliography}{9}
\bibitem{tangram_vision}
Tangram Vision Blog. \textit{Exploring Distortion and Distortion Models}. \\
Dostupné na: \url{https://www.tangramvision.com/blog/camera-modeling-exploring-distortion-and-distortion-models-part-iii}

\bibitem{opencv_colorspaces}
OpenCV Documentation. \textit{Changing Colorspaces}. \\
Dostupné na: \url{https://docs.opencv.org/4.x/df/d9d/tutorial_py_colorspaces.html}

\bibitem{opencv_track}
LearnOpenCV. \textit{Color Spaces in OpenCV}. \\
Dostupné na: \url{https://learnopencv.com/color-spaces-in-opencv-cpp-python/}

\bibitem{hsv_converter}
Online HSV Converter a kalkulačka pre OpenCV rozsahy. \\
Dostupné na: \url{https://vgyp.github.io/OpenCV-HSV-Color-Picker/}

\end{thebibliography}

\end{document}
